\chapter{文件系统实现}

\section{文件系统结构}
\concept{用户视图与系统视图}
\begin{points}
  \item 用户视图中，文件是带名字的逻辑数据集合，通过系统调用访问。
  \item 系统视图中，文件是磁盘数据块的集合，文件系统负责从逻辑文件偏移映射到物理磁盘块。
  \item 文件系统基本操作单位是数据块 block，而磁盘物理基本单位是扇区 sector。
  \item 操作系统通常一次读写一个块，块可由多个连续扇区组成。
\end{points}
\plain{同一个 `getc/putc` 看起来只动了 1 个字节，但系统层面往往还是先把整个数据块调入缓存，再在块内读写那个字节。}
\exam{判断题常考“文件系统的基本操作单位是块，不是字节也不是扇区”，这个层次关系要分清。}


\concept{分层结构}
\begin{points}
  \item 逻辑文件系统：处理文件名、目录、权限、FCB/inode 等。
  \item 文件组织模块：完成逻辑块到物理块的映射。
  \item 基本文件系统：向设备驱动发出读写块请求。
  \item I/O 控制层：设备驱动和中断处理程序控制硬件。
\end{points}
\plain{分层结构就是把“文件名解析、块映射、具体设备访问”拆开，各层各管一段。}
\exam{除职责分工外，PPT 还强调了代价：分层能降低复杂度和冗余，但会引入额外开销、降低性能。}


\concept{FCB 与 inode}
\begin{points}
  \item FCB 文件控制块保存文件属性和存储位置信息。
  \item inode 是类 Unix 系统中的索引节点，保存文件类型、权限、大小、时间、数据块地址等元数据。
  \item 目录项通常保存文件名和 inode 号，文件名不一定存放在 inode 内。
  \item inode 与文件物理分配方式有关，因为 inode 中需要保存指向数据块或索引块的指针。
\end{points}
\plain{FCB/inode 保存文件属性和物理位置，是文件系统找到文件内容的核心索引。}
\exam{常考目录项是否要放全部 FCB、inode 与文件名分离的好处。}

\concept{inode 号与按名查找过程}
\begin{points}
  \item 每个 inode 都有唯一的 \term{inode 号}，操作系统内部常用 inode 号而不是文件名来唯一标识文件。
  \item 目录项中通常保存“文件名 + inode 号”，因此目录的首要任务是把名字翻译成 inode 号。
  \item 按名查找文件时，系统先在目录中找到目标文件名对应的 inode 号，再据此读取 inode 信息，最后找到文件数据块。
  \item 因而文件名主要服务于用户使用，inode 号才是内核继续定位文件内容时真正依赖的内部标识。
\end{points}
\plain{可以把文件名看成给人看的“别名”，把 inode 号看成系统内部真正拿来追踪文件本体的“身份证号”。}
\exam{选择题常考“目录项里存的是什么”“系统内部按什么识别文件”“按名查找文件分几步完成”。}


\concept{盘上结构}
\begin{points}
  \item 引导控制块保存从该卷引导操作系统所需的信息。
  \item 卷控制块保存块总数、块大小、空闲块信息、空闲 FCB 或 inode 信息等。
  \item 目录结构负责组织文件名到 FCB/inode 的映射。
  \item FCB 或 inode 区保存每个文件的元数据和数据块地址。
\end{points}
\plain{盘上结构就是文件系统写在磁盘上的“管理骨架”，系统得先找到它们，后面才能找到文件。}
\exam{常考引导块、超级块、目录、inode 区分别干什么。}


\concept{内存结构}
\begin{points}
  \item 安装表 mount table 记录已挂载文件系统、挂载点和文件系统类型。
  \item 目录缓存保存最近访问过的目录信息，用于加速路径解析。
  \item 系统级打开文件表保存已打开文件的 FCB 副本、打开计数和磁盘位置等共享信息。
  \item 进程级打开文件表保存文件描述符、访问权限、当前读写位置等进程私有信息。
\end{points}
\plain{盘上结构负责“长期保存”，内存结构负责“当前加速访问”；其中 buffer 更偏传输缓冲，cache 更偏复用热点数据和元数据。}
\exam{PPT 明确区分了 buffer 和 cache，选择题里不要把两者完全当同义词。}


\concept{open 与 close 的实现}
\begin{points}
  \item open 先查系统打开文件表，若文件已被别的进程打开，则增加打开计数并在本进程表中建立指针。
  \item 若系统表中没有该文件，则搜索目录，把对应 FCB/inode 装入系统打开文件表，再返回文件句柄。
  \item close 删除进程级打开文件表项，并将系统级打开计数减一。
  \item 当打开计数降为 0 时，才把需要写回的元数据更新到磁盘并释放系统表项。
\end{points}
\plain{open 的真正价值是把后续读写反复要用的控制信息先搬进内存，不必每次都重新搜目录。}
\exam{常考 open 为什么能减少访盘，以及多个进程同时打开同一文件时哪些信息共享、哪些不共享。}


\section{目录实现}
\concept{线性列表目录}
\begin{points}
  \item 目录项按线性表保存，查找时顺序扫描。
  \item 优点：实现简单。
  \item 缺点：目录项多时查找慢，平均查找约需扫描一半目录项。
  \item 目录项访问次数题常用“目录项数/每块目录项数/平均扫描块数”计算。
\end{points}
\plain{线性目录查找慢的根因是要顺着比对名字，目录越大，平均扫描项数越多。}
\exam{常考线性目录为什么实现简单但查找代价高，以及平均查找长度估算。}


\concept{哈希目录}
\begin{points}
  \item 根据文件名计算哈希值，快速定位目录项。
  \item 优点：查找速度快。
  \item 缺点：需要处理哈希冲突，目录扩展和哈希表维护更复杂。
\end{points}
\plain{哈希目录的价值在于按名字快速定位，但冲突处理做不好，性能和管理复杂度都会出问题。}
\exam{常考哈希目录相对线性目录的优缺点，尤其冲突问题。}


\concept{inode 对目录检索的加速作用}
\begin{points}
  \item 若目录项中直接保存完整 FCB，则每个目录块能容纳的目录项数量较少。
  \item 把文件名和其他描述信息分离后，目录项只保留“文件名 + inode 号”，一个目录块能装下更多文件名。
  \item 检索时先按文件名找到 inode 号，再读取对应 inode。
  \item 因此按名查找目录时，平均访盘次数通常会下降。
\end{points}
\plain{inode 不是为了让结构更复杂，而是为了让“按文件名查找”这件事更省盘块访问。}
\exam{计算题常用平均访盘次数公式：平均扫描目录块数再加是否还要额外访问 inode 的次数，过程要写清。}


\section{文件存储空间分配方法}
\concept{文件分配的目标}
\begin{points}
  \item 文件存储空间分配首先要回答两个问题：如何为文件分配磁盘块，以及如何在高效利用磁盘空间的同时保证文件访问速度。
  \item 理想目标通常包括：提高磁盘空间利用率、支持文件动态增长、便于顺序访问与随机访问、减少外部碎片与管理开销。
  \item 不同分配方法本质上是在“访问性能、空间利用率、扩展方便性、实现复杂度”之间做权衡。
\end{points}
\plain{这一节不是一上来就背连续/链式/索引三种方法，而是先看它们在解决什么矛盾：既想省空间，又想访问快，还想文件好扩展。}
\exam{简答题若问“文件分配方法设计时要考虑什么”，可以按“空间利用率、访问效率、动态增长、管理开销”四点组织答案。}

\concept{连续分配}
\begin{points}
  \item 文件占用磁盘上一组连续块，目录项只需记录起始块号和长度。
  \item 优点：顺序访问速度快，随机访问也容易，寻道少。
  \item 缺点：需要预先知道文件大小，文件增长困难，会产生外部碎片。
  \item 适合大小较稳定、顺序访问多的文件。
\end{points}
\plain{连续分配答题抓三点：定位快、顺序访问好、增长困难且有外部碎片。}
\exam{常考连续分配的优缺点，以及为什么它不利于文件动态增长。}


\concept{链式分配}
\begin{points}
  \item 文件由多个磁盘块组成链表，每个块包含指向下一个块的指针。
  \item 优点：无外部碎片，文件容易增长。
  \item 缺点：随机访问差，访问第 $k$ 个块必须从第一个块顺链找；指针损坏会影响后续数据。
  \item FAT 是链式分配的改进，把链接指针集中存放在文件分配表中。
\end{points}
\plain{每个块指向下一个块，不怕文件增长，但找第 n 块必须顺着链走。}
\exam{PPT 还给了 FAT 表空间计算题：先算总块数，再算块号至少要多少位，最后乘每表项大小。}


\concept{FAT 表空间计算}
\begin{points}
  \item FAT 中通常“每个磁盘块/簇对应一个表项”，所以先用磁盘容量除以块大小，求出总块数。
  \item 再根据总块数判断“块号至少需要多少位”才能编号全部块。
  \item 为方便存取，实际表项长度常向上取到便于组织的字节数，如 12 位、16 位、20 位等。
  \item FAT 总空间 $=$ 表项大小 $\times$ 总块数。
\end{points}
\plain{FAT 计算题的套路很固定：先算有多少块，再算编号这些块至少要多少位，最后乘出整张 FAT 的大小。}
\exam{PPT 直接举了 `1.2MB/1KB` 和 `200MB/1KB` 这类例子，做题时不要跳过“块号要几位、表项实际取几位”这一步。}


\concept{索引分配}
\begin{points}
  \item 为每个文件建立索引块，索引块中保存该文件所有数据块号。
  \item 目录项保存索引块地址。
  \item 优点：支持直接访问，无外部碎片，文件可离散存放。
  \item 缺点：索引块本身占空间；小文件可能浪费索引块，大文件需要多级索引。
\end{points}
\plain{把文件所有数据块号集中放在索引块里，既支持离散分配，又方便随机访问。}
\exam{常考单级/多级索引最大文件长度计算：每个索引块能放多少块号，再乘数据块大小。}


\concept{多级索引与最大文件长度}
\begin{points}
  \item 若块大小为 $B$ 字节，每个块号占 $p$ 字节，则一个索引块可保存 $B/p$ 个块号。
  \item 一级索引最大可索引 $(B/p)$ 个数据块，最大文件大小为 $(B/p)\times B$。
  \item 二级索引最大可索引 $(B/p)^2$ 个数据块，最大文件大小为 $(B/p)^2\times B$。
  \item 三级索引同理为 $(B/p)^3\times B$。
\end{points}
\plain{这类题的套路固定：先算一个索引块能装多少块号，再根据几级索引逐层乘开，最后再乘数据块大小。}
\exam{不要漏掉“直接地址项”和“间接地址项”混合存在的情况，Unix/Linux 常考的是混合索引而不是纯多级索引。}


\concept{混合索引分配}
\begin{points}
  \item 混合索引把直接地址、一次间接、二次间接、三次间接结合起来，兼顾小文件效率与大文件容量。
  \item 典型 Unix inode 中常见 12 个直接地址项、1 个一次间接、1 个二次间接、1 个三次间接地址项。
  \item 小文件可直接用直接地址快速定位，大文件再逐级借助间接索引扩展容量。
  \item 计算最大文件长度时，要把各部分寻址范围分别算出后相加。
\end{points}
\plain{混合索引的设计思想是：多数文件其实不大，没必要一上来就走多级索引；但系统又必须给超大文件留扩展空间。}
\exam{考研题很爱出这种“若有若干直接地址、若干一级间接、一个二级间接”的最大文件长度计算。}

\section{空闲空间管理}
\concept{位示图/位图}
\begin{points}
  \item 用一个二进制位表示一个磁盘块是否空闲。
  \item 优点：查找连续空闲块方便，结构紧凑。
  \item 缺点：磁盘很大时位图本身也较大，需要常驻或缓存。
\end{points}
\plain{位图法的关键是“按位记录每个块是否空闲”，所以做题常要把字序号、位偏移和实际块号互相换算。}
\exam{PPT 给了大磁盘下位图本身占多大空间的估算题：先算总块数，再除以 8 得到字节数。}


\concept{空闲链表}
\begin{points}
  \item 把所有空闲块链接成链表。
  \item 优点：释放和分配单个块较简单。
  \item 缺点：查找连续空闲块效率低，链指针损坏可能影响管理。
\end{points}
\plain{链接法的主要短板不是“不能管空闲块”，而是“很难快速拿到一大片连续空闲块”。}
\exam{若题目强调连续空间分配需求，通常位图法或计数法更合适，单纯空闲链表往往不优。}


\concept{成组链接法}
\begin{points}
  \item 把空闲块分组，每组第一块记录下一组空闲块号。
  \item 兼顾链表和批量管理，是 Unix 早期常用方法。
  \item 适合大量空闲块的快速分配和回收。
\end{points}
\plain{成组链接法不是逐块找空闲块，而是把很多空闲块号成组保存，减少管理链过长的问题。}
\exam{PPT 对 Unix 成组链接法的“分配最后一个块号时先装入下一组”“表满回收时先整组写回”讲得很细，这两步是流程题重点。}


\concept{Unix 成组链接法的分配与回收}
\begin{points}
  \item 分配空闲块时，先从内存中当前专用块的空闲块号表取出下一个可用块号。
  \item 若取出的正好是当前组记录的“最后一个块号”，则要先把该块中保存的下一组空闲块号表读入内存，再把这个块分配出去。
  \item 回收空闲块时，若当前专用块中的空闲块号表未满，就直接把回收块号填入表中。
  \item 若空闲块号表已满，则先把当前整张空闲块号表写入这个新回收的块中，再把该回收块号作为新组首块号放回表中重新开始计数。
\end{points}
\plain{成组链接法流程题的关键是分清两种“临界时刻”：分配到当前组最后一个块号时要先装入下一组；回收时若表已满，要先把整组信息塞进新回收块。}
\exam{这类题通常不难算，难在顺序别写反；一旦把“先读下一组再分配”或“先写满表再重置新组”写错，后面全错。}


\concept{计数法}
\begin{points}
  \item 用“起始空闲块号 + 连续空闲块数”表示一段连续空闲区。
  \item 适合磁盘空闲块成片连续的情况。
  \item 比逐块链表更紧凑，但碎片多时记录数量增加。
\end{points}
\plain{计数法适合连续空闲区较多的磁盘，用一条记录概括一整段连续空闲块。}
\exam{判断题会考它为什么适合“频繁、连续的分配与释放”，核心原因是连续空闲段能被压缩成一条记录。}


\section{文件共享实现}
\concept{文件共享的动机}
\begin{points}
  \item 文件共享允许不同用户共同使用同一个文件。
  \item 这样做的主要动机有三类：支持用户合作、减少磁盘空间开销、减少同一数据维护多份副本带来的不一致性。
  \item 因而文件共享不仅是“方便别人读我的文件”，更是从协作效率、存储成本和一致性控制三方面做优化。
\end{points}
\plain{共享的价值不只是省空间；更关键的是多人协作时大家面对同一份数据，能少掉很多“各改各的副本然后版本打架”的问题。}
\exam{若选择题问“为什么要实现文件共享”，标准关键词就是合作、节省空间、减少不一致。}


\concept{基于索引节点的共享}
\begin{points}
  \item 多个目录项指向同一个 inode，实现硬链接共享。
  \item inode 中维护引用计数，记录有多少目录项引用它。
  \item 删除目录项只减少引用计数，不一定删除文件内容。
\end{points}
\plain{硬链接共享的是同一个 inode，所以任何一个用户修改文件内容，其他硬链接看到的都是同一份数据。}
\exam{还要会写出结论：创建硬链接时引用计数加一；只有所有硬链接都删掉后，文件才真正删除。}


\concept{索引节点中的链接计数}
\begin{points}
  \item inode 中的链接计数字段 \code{count} 表示当前有多少个目录项指向该 inode。
  \item 创建硬链接时，\code{count} 加 1；删除一个硬链接时，\code{count} 减 1。
  \item 只要 \code{count} 仍大于 0，文件数据通常不会真正删除，因为仍有其他文件名可访问它。
  \item 因而“删掉某个文件名”与“文件内容真正消失”在硬链接语义下不是同一件事。
\end{points}
\plain{硬链接题的关键是把“名字”和“文件本体”分开：目录项删掉的首先只是一个名字，只有最后一个名字也删掉时，文件才真正失去引用。}
\exam{PPT 单独把 \code{count} 拿出来讲，选择题和简答题都常考：什么时候加一、减一，以及文件真正删除的条件。}


\concept{基于符号链接的共享}
\begin{points}
  \item 符号链接文件保存目标路径名。
  \item 访问符号链接时，系统根据路径再去查找目标文件。
  \item 优点是灵活、可跨文件系统；缺点是访问多一次路径解析，目标删除后链接失效。
\end{points}
\plain{符号链接并不直接共享 inode，它只是保存“去哪里找目标文件”的路径，因此本身更像快捷方式。}
\exam{常考硬链接与软链接对引用计数的影响：硬链接会增加原文件引用计数，软链接通常不会。}
